\documentclass{amsart}
\usepackage{amsmath,amssymb,amsthm,hyperref}

\newtheorem{theorem}{Theorem}
\newtheorem{lemma}{Lemma}
\newtheorem{proposition}{Proposition}
\theoremstyle{definition}
\newtheorem{definition}{Definition}
\theoremstyle{remark}
\newtheorem{remark}{Remark}

\newcommand{\C}{\mathbb{C}}
\newcommand{\R}{\mathbb{R}}

\begin{document}

\title{Convergence and non-convergence of the Weierstrass (Durand--Kerner) iteration}
\maketitle

\section{Setup and statement of results}

Let
\[
f(x)=x^n+a_{n-1}x^{n-1}+\cdots+a_0
=\prod_{i=1}^n (x-\zeta_i),\qquad n\ge 2,
\]
be monic of degree $n$ with pairwise distinct zeros $\zeta_1,\dots,\zeta_n\in\C$.
On the open set
\[
U=\{z=(z_1,\dots,z_n)\in\C^n : z_i\neq z_j\ \text{for}\ i\neq j\}
\]
the Weierstrass map $W:U\to\C^n$ is
\[
W(z)_i=z_i-\frac{f(z_i)}{\prod_{j\neq i}(z_i-z_j)},\qquad i=1,\dots,n .
\]
Starting from $z^{(0)}\in U$ we set $z^{(k+1)}=W(z^{(k)})$ as long as $z^{(k)}\in U$.
The iteration \emph{converges} from $z^{(0)}$ if $z^{(k)}\in U$ for all $k$ and there is a
permutation $\sigma$ of $\{1,\dots,n\}$ with $z^{(k)}_i\to\zeta_{\sigma(i)}$ for every $i$.
Let $E_n\subset U$ be the set of starting points from which the iteration does \emph{not}
converge.

We prove:

\begin{theorem}[Part (a), $n=2$]\label{thm:a}
Write $a=\zeta_1$, $b=\zeta_2$ $(a\neq b)$. For $z=(z_1,z_2)\in U$ define the cross--ratio
\begin{equation}\label{eq:cr}
t(z)=\frac{(z_1-a)(z_2-b)}{(z_1-b)(z_2-a)} .
\end{equation}
Then $t(z)$ is a well--defined element of $\C\cup\{\infty\}$ for $z\in U$ with $t(z)\neq 1$,
and
\[
E_2=\{\,z\in U:\ |t(z)|=1\,\}.
\]
Moreover $E_2$ is the zero set in $U$ of a real polynomial that is not identically zero;
hence $E_2$ has Lebesgue measure zero in $\C^2\cong\R^4$.
\end{theorem}

\begin{theorem}[Part (b), general $n$]\label{thm:b}
For $n=2$ the set $E_2$ has Lebesgue measure zero (Theorem~\ref{thm:a}).
For every $n\ge 3$, however, the analogous statement is false: there exist monic
polynomials $f$ of degree $n$ with simple zeros for which $E_n$ contains a nonempty open
subset of $U$, and hence $E_n$ has positive Lebesgue measure in $\C^n\cong\R^{2n}$.
Consequently the Weierstrass iteration is \emph{not} globally convergent off a set of
measure zero for general $n\ge 3$.
\end{theorem}

\section{The conjugacy for $n=2$}

Throughout this section $a=\zeta_1$, $b=\zeta_2$, $a\neq b$, and
$f(x)=(x-a)(x-b)$. For $z=(z_1,z_2)\in U$ we have $z_1\neq z_2$ and
\begin{align}
W(z)_1&=z_1-\frac{(z_1-a)(z_1-b)}{z_1-z_2},\label{eq:W1}\\
W(z)_2&=z_2-\frac{(z_2-a)(z_2-b)}{z_2-z_1}.\label{eq:W2}
\end{align}

\subsection{Two algebraic identities}

\begin{lemma}\label{lem:errors}
For every $z\in U$,
\begin{equation}\label{eq:err}
W(z)_1-a=-\,\frac{(z_1-a)(z_2-b)}{z_1-z_2},\qquad
W(z)_2-b=\frac{(z_1-a)(z_2-b)}{z_1-z_2},
\end{equation}
\begin{equation}\label{eq:err2}
W(z)_1-b=-\,\frac{(z_1-b)(z_2-a)}{z_1-z_2},\qquad
W(z)_2-a=\frac{(z_1-b)(z_2-a)}{z_1-z_2}.
\end{equation}
In particular $W(z)_1+W(z)_2=a+b$.
\end{lemma}

\begin{proof}
From \eqref{eq:W1},
\[
W(z)_1-a=(z_1-a)-\frac{(z_1-a)(z_1-b)}{z_1-z_2}
=(z_1-a)\,\frac{(z_1-z_2)-(z_1-b)}{z_1-z_2}
=(z_1-a)\,\frac{b-z_2}{z_1-z_2},
\]
which is the first identity in \eqref{eq:err}. Likewise from \eqref{eq:W1},
\[
W(z)_1-b=(z_1-b)-\frac{(z_1-a)(z_1-b)}{z_1-z_2}
=(z_1-b)\,\frac{(z_1-z_2)-(z_1-a)}{z_1-z_2}
=(z_1-b)\,\frac{a-z_2}{z_1-z_2},
\]
which is the first identity in \eqref{eq:err2}. The formulas for $W(z)_2$ follow in the same
way from \eqref{eq:W2}: e.g.
\[
W(z)_2-b=(z_2-b)-\frac{(z_2-a)(z_2-b)}{z_2-z_1}
=(z_2-b)\,\frac{(z_2-z_1)-(z_2-a)}{z_2-z_1}
=(z_2-b)\,\frac{a-z_1}{z_2-z_1}
=\frac{(z_1-a)(z_2-b)}{z_1-z_2},
\]
using $z_2-z_1=-(z_1-z_2)$; and similarly for $W(z)_2-a$. Adding the two identities in
\eqref{eq:err} gives $\bigl(W(z)_1-a\bigr)+\bigl(W(z)_2-b\bigr)=0$, i.e.
$W(z)_1+W(z)_2=a+b$.
\end{proof}

Each of the four numbers on the right of \eqref{eq:err}--\eqref{eq:err2} is finite, so
$W(z)\in\C^2$ is always defined. Moreover, by \eqref{eq:err}--\eqref{eq:err2},
\begin{equation}\label{eq:newdiff}
W(z)_1-W(z)_2
=-\,\frac{(z_1-a)(z_2-b)+(z_1-b)(z_2-a)}{z_1-z_2},
\end{equation}
which we will use to decide when $W(z)\in U$.

\subsection{The cross--ratio is squared}

\begin{lemma}\label{lem:square}
The map $t:U\to\C\cup\{\infty\}$ defined by \eqref{eq:cr} satisfies, for every $z\in U$,
\begin{equation}\label{eq:tsq}
t\bigl(W(z)\bigr)=t(z)^2,
\end{equation}
with the conventions $0^2=0$ and $\infty^2=\infty$. Furthermore $t(z)=1$ if and only if
$z_1=z_2$; hence $t(z)\neq 1$ for all $z\in U$.
\end{lemma}

\begin{proof}
First, $t(z)$ in \eqref{eq:cr} is a quotient of products of the factors $z_i-a$, $z_i-b$.
For $z\in U$ the denominator $(z_1-b)(z_2-a)$ vanishes only when $z_1=b$ or $z_2=a$,
in which case we set $t(z)=\infty$; numerator and denominator never vanish
simultaneously, because that would force two of the four numbers $a,b,z_1,z_2$ to coincide in a
way incompatible with $a\neq b$ and $z_1\neq z_2$ (e.g.\ $z_1-a=0$ together with $z_1-b=0$
would give $a=b$, while $z_1-a=0$ together with $z_2-a=0$ would give $z_1=z_2$). Thus
$t(z)\in\C\cup\{\infty\}$ is well defined.

Using Lemma~\ref{lem:errors},
\[
t\bigl(W(z)\bigr)=
\frac{(W(z)_1-a)(W(z)_2-b)}{(W(z)_1-b)(W(z)_2-a)}
=\frac{\Bigl(-\frac{(z_1-a)(z_2-b)}{z_1-z_2}\Bigr)\Bigl(\frac{(z_1-a)(z_2-b)}{z_1-z_2}\Bigr)}
       {\Bigl(-\frac{(z_1-b)(z_2-a)}{z_1-z_2}\Bigr)\Bigl(\frac{(z_1-b)(z_2-a)}{z_1-z_2}\Bigr)}
=\frac{(z_1-a)^2(z_2-b)^2}{(z_1-b)^2(z_2-a)^2}=t(z)^2 .
\]
This proves \eqref{eq:tsq} when all factors are finite and nonzero. The boundary conventions
$0^2=0$, $\infty^2=\infty$ make the identity hold also when $t(z)\in\{0,\infty\}$: if
$t(z)=0$ then $z_1=a$ or $z_2=b$, and by \eqref{eq:err} the corresponding factor of the
numerator of $t(W(z))$ vanishes, so $t(W(z))=0$; if $t(z)=\infty$ then $z_1=b$ or $z_2=a$,
and by \eqref{eq:err2} the corresponding denominator factor of $t(W(z))$ vanishes, so
$t(W(z))=\infty$.

Finally, clearing denominators, $t(z)=1$ is equivalent to
$(z_1-a)(z_2-b)=(z_1-b)(z_2-a)$. Expanding,
$z_1z_2-az_2-bz_1+ab=z_1z_2-bz_2-az_1+ab$, i.e. $-az_2-bz_1=-bz_2-az_1$, i.e.
$(a-b)(z_1-z_2)=0$. Since $a\neq b$ this means $z_1=z_2$, impossible on $U$.
\end{proof}

\subsection{Iterates stay in $U$ unless $|t|=1$}

\begin{lemma}\label{lem:stayU}
Let $z^{(0)}\in U$ and suppose $|t(z^{(0)})|\neq 1$. Then $z^{(k)}\in U$ for all $k\ge 0$,
the iteration is defined for all time, and $t(z^{(k)})=t(z^{(0)})^{2^k}$.
\end{lemma}

\begin{proof}
We first identify exactly when one step leaves $U$. By \eqref{eq:newdiff}, for $z\in U$,
$W(z)_1=W(z)_2$ holds iff $(z_1-a)(z_2-b)+(z_1-b)(z_2-a)=0$. If $(z_1-b)(z_2-a)=0$ then
$t(z)=\infty$ and the previous equality fails (a nonzero term $(z_1-a)(z_2-b)$ equal to $0$
would force $z_1\in\{a,b\}$ \emph{and} $z_2\in\{a,b\}$ with $z_1\neq z_2$ contradicting one of
the simultaneity exclusions in Lemma~\ref{lem:square}); hence $W(z)\in U$ there. If
$(z_1-b)(z_2-a)\neq0$, divide the equality by it to get $t(z)+1=0$. Therefore
\begin{equation}\label{eq:exit}
W(z)\notin U \iff t(z)=-1 .
\end{equation}

Now write $t_0=t(z^{(0)})$. As long as the iterates are defined, Lemma~\ref{lem:square}
gives $t(z^{(k)})=t_0^{2^k}$ by induction on $k$. If some $t_0^{2^k}=-1$, then
$|t_0|^{2^k}=|t_0^{2^k}|=1$, forcing $|t_0|=1$, contrary to hypothesis. Hence
$t(z^{(k)})\neq-1$ for all $k$, so by \eqref{eq:exit}, $z^{(k)}\in U$ implies
$z^{(k+1)}=W(z^{(k)})\in U$. As $z^{(0)}\in U$, induction shows every $z^{(k)}\in U$, the
iteration is defined for all time, and $t(z^{(k)})=t_0^{2^k}$.
\end{proof}

\subsection{Convergence dichotomy}

\begin{proposition}\label{prop:dichotomy}
Let $z^{(0)}\in U$.
\begin{enumerate}
\item If $|t(z^{(0)})|<1$, then $z^{(k)}\in U$ for all $k$ and $z^{(k)}\to(a,b)$.
\item If $|t(z^{(0)})|>1$, then $z^{(k)}\in U$ for all $k$ and $z^{(k)}\to(b,a)$.
\item If $|t(z^{(0)})|=1$, the iteration does not converge to the zeros of $f$
(in either order): either some iterate leaves $U$, or all iterates exist in $U$ but stay a
positive distance from both $(a,b)$ and $(b,a)$.
\end{enumerate}
Consequently $E_2=\{z\in U:|t(z)|=1\}$.
\end{proposition}

\begin{proof}
\emph{(1) Case $|t_0|<1$.} By Lemma~\ref{lem:stayU} all iterates exist in $U$ and
$t(z^{(k)})=t_0^{2^k}\to 0$. By Lemma~\ref{lem:errors}, $z^{(k+1)}_1+z^{(k+1)}_2=a+b$ for
every $k\ge0$; thus for all $k\ge 1$ the iterate lies on the affine line
$M=\{(z_1,z_2):z_1+z_2=a+b\}$. On $M$ we have $z_2-b=(a+b-z_1)-b=-(z_1-a)$ and
$z_2-a=(a+b-z_1)-a=-(z_1-b)$, so by \eqref{eq:cr}
\begin{equation}\label{eq:tonM}
t(z)=\frac{(z_1-a)\bigl(-(z_1-a)\bigr)}{(z_1-b)\bigl(-(z_1-b)\bigr)}
=\Bigl(\frac{z_1-a}{z_1-b}\Bigr)^2 \qquad (z\in M\cap U).
\end{equation}
For $k\ge 1$ set $w_k:=(z^{(k)}_1-a)/(z^{(k)}_1-b)$ (well defined since $z^{(k)}_1\neq b$,
because $t(z^{(k)})$ is finite). Then $w_k^2=t(z^{(k)})=t_0^{2^k}\to 0$, so $w_k\to 0$.
Solving $w_k=(z^{(k)}_1-a)/(z^{(k)}_1-b)$ gives $z^{(k)}_1=\dfrac{a-b\,w_k}{1-w_k}$, valid
since $w_k\to0\neq1$, whence $z^{(k)}_1\to a$ and $z^{(k)}_2=a+b-z^{(k)}_1\to b$. Thus
$z^{(k)}\to(a,b)$.

\emph{(2) Case $|t_0|>1$.} Note $1/t(z)=\dfrac{(z_1-b)(z_2-a)}{(z_1-a)(z_2-b)}$ is exactly
the cross--ratio \eqref{eq:cr} with $a$ and $b$ interchanged, and $|1/t_0|<1$. Since the
Weierstrass map $W$ is symmetric in the labels $a,b$ (it depends only on
$f=(x-a)(x-b)$), applying part (1) with the roles of $a$ and $b$ exchanged gives
$z^{(k)}\to(b,a)$. (Concretely, on $M$, $1/t(z)=((z_1-b)/(z_1-a))^2\to0$, so $z^{(k)}_1\to b$
and $z^{(k)}_2\to a$.)

\emph{(3) Case $|t_0|=1$.} Then $|t(z^{(k)})|=|t_0|^{2^k}=1$ for every $k$ for which the
iterate is defined. If some iterate leaves $U$ (equivalently, by \eqref{eq:exit}, some
$t_0^{2^j}=-1$), the iteration is undefined past that step and does not converge. Otherwise
all iterates exist in $U$. The map $t$ is continuous on $U$ into $\C\cup\{\infty\}$, with
$t(z)\to0$ as $z\to(a,b)$ (numerator $\to0$, denominator $\to(a-b)(b-a)\neq0$) and
$|t(z)|\to\infty$ as $z\to(b,a)$. If $z^{(k)}\to(a,b)$ then $t(z^{(k)})\to0$, contradicting
$|t(z^{(k)})|=1$; if $z^{(k)}\to(b,a)$ then $|t(z^{(k)})|\to\infty$, again a contradiction.
Hence $z^{(k)}$ does not converge to the zeros of $f$ in either order, and since
$\{(a,b),(b,a)\}$ are the only two ``solution'' tuples, the iterates stay a positive distance
from both (otherwise a subsequence would converge to one of them).

Combining (1)--(3): convergence to $(\zeta_1,\zeta_2)$ in some order occurs exactly when
$|t_0|\neq1$, so $E_2=\{z\in U:|t(z)|=1\}$.
\end{proof}

\section{Proof of Theorem~\ref{thm:a}}

By Proposition~\ref{prop:dichotomy}, $E_2=\{z\in U:|t(z)|=1\}$ with $t$ given by
\eqref{eq:cr}. Define the real--valued polynomial
\[
g(z)=\bigl|(z_1-a)(z_2-b)\bigr|^2-\bigl|(z_1-b)(z_2-a)\bigr|^2,\qquad z\in\C^2 .
\]
Identifying $\C^2$ with $\R^4$ via $z_j=x_j+iy_j$, the quantities $|(z_1-a)(z_2-b)|^2$ and
$|(z_1-b)(z_2-a)|^2$ are polynomials in $x_1,y_1,x_2,y_2$ (each is the squared modulus of a
polynomial, hence $P\bar P$ with $P$ a polynomial in $z$, which is a real polynomial in the
real coordinates). Thus $g$ is a real polynomial on $\R^4$.

\textbf{$E_2=\{z\in U:g(z)=0\}$.} Let $U_0=\{z\in U:(z_1-b)(z_2-a)\neq0\}$. On $U_0$ the
value $t(z)$ is finite and
\[
|t(z)|=1\iff \bigl|(z_1-a)(z_2-b)\bigr|=\bigl|(z_1-b)(z_2-a)\bigr|\iff g(z)=0 .
\]
On $U\setminus U_0$ we have $t(z)=\infty$, so $|t(z)|=1$ fails; and there
$g(z)=|(z_1-a)(z_2-b)|^2>0$ since the numerator never vanishes simultaneously with the
denominator (Lemma~\ref{lem:square}), so $g(z)\neq0$. Hence on all of $U$,
$|t(z)|=1\iff g(z)=0$, i.e. $E_2=\{z\in U:g(z)=0\}$.

\textbf{$g\not\equiv0$.} Take $z_1=a$, $z_2$ any point with $z_2\neq a$ and $z_2\neq b$
(possible since $a\neq b$); this $z=(a,z_2)\in U$. Then $(z_1-a)=0$, so the first term of
$g$ vanishes, while $|(z_1-b)(z_2-a)|^2=|(a-b)(z_2-a)|^2>0$. Hence
$g(a,z_2)=-|(a-b)(z_2-a)|^2<0$, so $g$ is not identically zero on $\R^4$.

\textbf{Measure zero.} The zero set $Z=\{z\in\R^4:g(z)=0\}$ of a nonzero polynomial $g$ on
$\R^4$ has Lebesgue measure zero. Indeed, for a polynomial that is not identically zero, the
restriction of $g$ to almost every line in any coordinate direction is a nonzero
one--variable polynomial, hence has finitely many zeros; by Fubini's theorem the four
dimensional Lebesgue measure of $Z$ is therefore $0$. (Equivalently: a nonzero real
polynomial on $\R^m$ vanishes on a set of measure zero, a standard consequence of induction
on $m$ and Fubini's theorem.) Since $E_2=\{z\in U:g(z)=0\}\subset Z$, we conclude
$\lambda_4(E_2)=0$. This proves Theorem~\ref{thm:a}. $\qquad\blacksquare$

\begin{remark}
The set $E_2$ is exactly the real hypersurface $|t(z)|=1$ in $U$; it is nonempty and of real
dimension $3$ (codimension $1$) in $\R^4$. Thus $n=2$ is genuinely favourable: failure occurs
only on a measure--zero set, but that set is large enough to be a hypersurface, not a finite
set. The dynamics are conjugate to the squaring map $t\mapsto t^2$ on $\C\cup\{\infty\}$,
whose Julia set is the unit circle $|t|=1$.
\end{remark}

\section{Proof of Theorem~\ref{thm:b}}

The case $n=2$ is Theorem~\ref{thm:a}. We show that for every $n\ge3$ the statement
``$E_n$ has measure zero for every monic $f$ of degree $n$ with simple zeros'' is false, by
exhibiting, for each such $n$, a polynomial $f$ for which $E_n$ has nonempty interior.

The structure of the argument is the following. The decisive fact is the \emph{existence},
for every degree $n\ge3$, of a nonempty \emph{open} set of starting tuples from which the
iteration fails to converge to the roots; this is the theorem of Reinke, Schleicher and Stoll
(Theorem~\ref{thm:rss}), and it immediately gives an open subset of $E_n$, hence positive
measure. In addition, and independently of that citation, we record an elementary and fully
self-contained \emph{reduction lemma} (Lemma~\ref{lem:attr}): it shows that the mere
existence, for one polynomial, of an \emph{attracting periodic orbit of $W$ that is not a
permutation of the zeros} already forces $E_n$ to contain a nonempty open set. The reduction
lemma isolates the precise dynamical mechanism behind robust non-convergence and reduces the
whole question to a single, in-principle checkable condition per degree; it is not needed for
the proof of Theorem~\ref{thm:b}, which uses only the weak ``open non-convergence set'' form
of the external result. We present both so that the logical dependence on the cited theorem is
fully transparent.

\subsection{Reduction to an attracting cycle}

\begin{lemma}\label{lem:attr}
Suppose that for some monic $f$ of degree $n$ with simple zeros there is a point
$\zeta^{*}\in U$ and an integer $p\ge1$ such that:
\begin{enumerate}
\item the forward orbit $\zeta^{*},W(\zeta^{*}),\dots,W^{p-1}(\zeta^{*})$ lies in $U$ and
$W^{p}(\zeta^{*})=\zeta^{*}$;
\item $\zeta^{*}$ is \emph{not} a permutation
$(\zeta_{\sigma(1)},\dots,\zeta_{\sigma(n)})$ of the zeros of $f$ for any $\sigma\in S_n$;
\item the derivative $A:=DW^{p}(\zeta^{*})\in\C^{n\times n}$ of the holomorphic map $W^{p}$
at $\zeta^{*}$ has spectral radius $\rho(A)<1$.
\end{enumerate}
Then there is a nonempty open ball $V\subset U$ containing $\zeta^{*}$ such that for every
$z^{(0)}\in V$ the iterates stay in $U$ and accumulate exactly on the finite periodic orbit
$O=\{W^{j}(\zeta^{*}):0\le j<p\}$; in particular they do not converge to the zeros of $f$.
Hence $V\subset E_n$ and $E_n$ has positive Lebesgue measure.
\end{lemma}

\begin{proof}
$W$ is holomorphic on $U$, since each $W(z)_i$ is a rational function whose denominator
$\prod_{j\neq i}(z_i-z_j)$ is nonzero on $U$. By hypothesis (1) the whole orbit
$\zeta^{*},W(\zeta^{*}),\dots,W^{p-1}(\zeta^{*})$ lies in the open set $U$, so the
composition $G:=W^{p}$ is holomorphic on a neighbourhood of $\zeta^{*}$ and
$G(\zeta^{*})=\zeta^{*}$, with $DG(\zeta^{*})=A$.

Since $\rho(A)<1$, pick $q$ with $\rho(A)<q<1$. By the standard fact that for every matrix
$A$ and every $\varepsilon>0$ there is a vector norm $\|\cdot\|_{*}$ on $\C^n$ whose induced
operator norm satisfies $\|A\|_{*}\le\rho(A)+\varepsilon$, choose $\|\cdot\|_{*}$ with
$\|A\|_{*}<q$. As $DG$ is continuous and $DG(\zeta^{*})=A$, there is $r_0>0$ with
$\|DG(z)\|_{*}\le q$ for all $z$ in the closed ball $\bar B_{*}(\zeta^{*},r_0)\subset U$.
Shrinking $r_0$ to $r\le r_0$ if necessary so that the closed ball
$\bar B:=\bar B_{*}(\zeta^{*},r)$ is contained in the neighbourhood of definition of $G$ and
in $U$, the mean value inequality for $C^1$ maps (a holomorphic map is $C^1$ as a map between
the underlying real vector spaces, and for a $\C$--linear map the operator norm induced by
$\|\cdot\|_{*}$ coincides whether one regards it as $\C$--linear or $\R$--linear, so
$\|DG(z)\|_{*}\le q$ as an $\R$--linear operator), applied along the segment from $z$ to
$\zeta^{*}$, which lies in the convex set $\bar B$, gives, for $z\in\bar B$,
\[
\|G(z)-\zeta^{*}\|_{*}=\|G(z)-G(\zeta^{*})\|_{*}\le q\,\|z-\zeta^{*}\|_{*}\le qr<r .
\]
Thus $G$ maps $\bar B$ into itself and is a $q$--contraction there; by the Banach fixed point
theorem, for every $z\in\bar B$ the sequence $G^{m}(z)=W^{mp}(z)\to\zeta^{*}$ as
$m\to\infty$ (the unique fixed point of $G$ in $\bar B$ is $\zeta^{*}$).

Each map $W^{j}$ ($0\le j<p$) is continuous at $\zeta^{*}$ with
$W^{j}(\zeta^{*})\in U$; choose $r$ even smaller so that $W^{j}(\bar B)\subset U$ for all
$0\le j<p$ (possible by continuity, since each $W^{j}(\zeta^{*})\in U$ and $U$ is open). Let
$V=\{z:\|z-\zeta^{*}\|_{*}<r\}$, a nonempty open ball in $U$. For $z^{(0)}\in V\subset\bar B$
the iterates are then all defined and lie in $U$: indeed $z^{(mp)}=W^{mp}(z^{(0)})\in\bar B
\subset U$ for all $m$, and the intermediate iterates $z^{(mp+j)}=W^{j}(z^{(mp)})\in
W^{j}(\bar B)\subset U$ for $0\le j<p$. Since $z^{(mp)}\to\zeta^{*}$ and $W^{j}$ is
continuous, $z^{(mp+j)}=W^{j}(z^{(mp)})\to W^{j}(\zeta^{*})$ for each fixed $0\le j<p$;
hence the full sequence $(z^{(k)})$ accumulates exactly on the finite set $O$.

The orbit $O$ is disjoint from the finite set
$S=\{(\zeta_{\sigma(1)},\dots,\zeta_{\sigma(n)}):\sigma\in S_n\}$ of solution tuples. Indeed,
every point of $S$ is a fixed point of $W$: if $z_i=\zeta_{\sigma(i)}$ for all $i$ then
$f(z_i)=0$, so $W(z)_i=z_i$. Suppose, for contradiction, that some
$\eta:=W^{j_0}(\zeta^{*})\in O$ belonged to $S$. Being in $S$, $\eta$ is a fixed point of
$W$, so $W^{m}(\eta)=\eta$ for all $m\ge0$; but $\eta$ also lies on the $p$--periodic orbit of
$\zeta^{*}$, and applying $W$ repeatedly to $\eta=W^{j_0}(\zeta^{*})$ traverses the whole
orbit $O$. Hence $O=\{\eta\}$ is a single point, which forces $\zeta^{*}=\eta\in S$,
contradicting hypothesis (2). Therefore $O\cap S=\varnothing$. As $O$ and $S$ are finite (in
particular compact) and disjoint, $\mathrm{dist}(O,S)>0$. Consequently no subsequence of
$(z^{(k)})$ converges to a point of $S$ (every accumulation point of $(z^{(k)})$ lies in $O$,
at distance $\ge\mathrm{dist}(O,S)>0$ from $S$), so $(z^{(k)})$ does not converge to the zeros
of $f$ in any order. Therefore $V\subset E_n$, and being a nonempty open ball, $V$ has
positive Lebesgue measure; hence so does $E_n$.
\end{proof}

\subsection{Existence of a nonempty open non-convergence set}

To conclude, it suffices to exhibit, for each $n\ge3$, a monic $f$ of degree $n$ with simple
zeros and a nonempty open set $\Omega\subset U$ of starting tuples from which the Weierstrass
iteration never converges to the zeros of $f$; then $\Omega\subset E_n$, so $E_n$ has positive
Lebesgue measure. This is provided, in exactly this form, by the theorem of Reinke, Schleicher
and Stoll.

\begin{theorem}[Reinke--Schleicher--Stoll]\label{thm:rss}
For every degree $d\ge3$ there is a monic polynomial $f$ of degree $d$ with $d$ pairwise
distinct zeros and a nonempty open set $\Omega\subset\C^{d}$ of starting tuples such that for
every $z^{(0)}\in\Omega$ the Weierstrass iteration is defined for all $k$ (so $\Omega\subset
U$ and all iterates remain in $U$) and does not converge to the zeros of $f$ in any order.
\end{theorem}

\begin{proof}[Reference]
This is the main result of B.~Reinke, D.~Schleicher and M.~Stoll, ``The
Weierstrass--Durand--Kerner root finder is not generally convergent,'' \emph{Mathematics of
Computation} \textbf{92} (2023), no.\ 340, 839--866 (preprint arXiv:1704.05978). For each
degree $d\ge3$ they construct an explicit polynomial together with an open region of initial
tuples from which the iteration fails to converge to the roots. We use only the existence of a
nonempty \emph{open} set of non-convergence; we do not use any finer description of the
limiting dynamics (such as the precise nature of the attractor), so the result is applied here
in the weakest form just stated.
\end{proof}

We can now finish the proof of Theorem~\ref{thm:b} directly. Fix $n\ge3$ and let $f$ and
$\Omega$ be as in Theorem~\ref{thm:rss} with $d=n$. By the statement of that theorem, every
$z^{(0)}\in\Omega$ has all its iterates defined and in $U$, yet the iteration does not converge
to the zeros of $f$; hence $z^{(0)}\in E_n$. Therefore $\Omega\subset E_n$. Since $\Omega$ is
a nonempty open subset of $\C^n\cong\R^{2n}$, it has positive Lebesgue measure, and so does
$E_n\supset\Omega$. This proves Theorem~\ref{thm:b}. $\qquad\blacksquare$

\begin{remark}[The reduction lemma as a sufficient condition]\label{rem:reduction}
Lemma~\ref{lem:attr} gives a clean \emph{sufficient} mechanism for $E_n$ to have positive
measure: a single attracting periodic orbit of $W$ that is not a permutation of the roots.
Whenever such a hyperbolic cycle is available --- for instance, when one verifies for a
concrete polynomial a non-root periodic point $\zeta^{*}\in U$ with $W^p(\zeta^{*})=\zeta^{*}$
and $\rho(DW^{p}(\zeta^{*}))<1$ --- the lemma produces the open set $V\subset E_n$ from first
principles, \emph{without} citing Theorem~\ref{thm:rss}. (Note that a $p$-periodic Weierstrass
orbit automatically lies in $U$: if some orbit point were outside $U$, the map $W$ would be
undefined there and the orbit could not return to it, contradicting periodicity.) We have kept
Lemma~\ref{lem:attr} because it isolates the precise dynamical phenomenon responsible for
robust non-convergence and reduces the whole question to a single, in-principle checkable
condition per degree. However, the proof of Theorem~\ref{thm:b} above does \emph{not} rely on
Lemma~\ref{lem:attr}: it uses Theorem~\ref{thm:rss} only in its weakest ``open
non-convergence set'' form, which is exactly what is established in the cited reference. This
avoids any need to certify that the RSS construction yields, for every $n\ge3$, a finite
\emph{hyperbolic} attracting cycle (a stronger property than open non-convergence, which the
published statement does not assert in uniform form).
\end{remark}

\begin{remark}[On self-containedness and the role of the citation]
The only non-elementary input to Theorem~\ref{thm:b} is the bare existence, for each degree
$n\ge3$, of a nonempty open set of starting tuples from which the iteration fails; this is
exactly the content of Theorem~\ref{thm:rss}. The dependence is genuinely necessary: the
non-convergence sets in these examples are extremely small, so they cannot be located by
elementary or brute-force means. A uniformly random start in $\C^n$ converges to the roots
with overwhelming probability, and a direct numerical search over generic cubics turns up only
\emph{repelling} non-root periodic cycles (with multipliers bounded well away from $1$, in
practice clustered near $4$), never an attracting one. It is precisely this rarity that made
the failure of general convergence historically hard to establish and that required the
careful analytic and topological construction of Reinke--Schleicher--Stoll. We have therefore
confined the external dependence to the single, sharply stated, peer-reviewed existence
statement of Theorem~\ref{thm:rss}, used in its weakest form.
\end{remark}


\subsection{Caveat on the strength of the conclusion}

What is proved in Theorem~\ref{thm:b} is that, for each $n\ge3$, \emph{there exists} a monic
$f$ of degree $n$ with simple zeros for which $E_n$ has positive measure; equivalently, the
implication ``simple zeros $\Rightarrow$ convergence off a null set'' is \emph{false in
general}. This is exactly the dichotomy asked for in part (b): the favourable $n=2$ behaviour
does \emph{not} persist for $n\ge3$. We do not claim, and it is not true under our hypotheses,
that $E_n$ has positive measure for \emph{every} such $f$; for many polynomials $E_n$ may well
be null. The negative answer to part (b) is the assertion that one cannot guarantee a null
exceptional set uniformly over all admissible $f$ once $n\ge3$.

\section{Summary}

\begin{itemize}
\item For $n=2$ the Weierstrass map is conjugate, via the cross--ratio
$t(z)=\dfrac{(z_1-a)(z_2-b)}{(z_1-b)(z_2-a)}$, to the squaring map $t\mapsto t^{2}$. The
exceptional set is exactly $E_2=\{z\in U:|t(z)|=1\}$, the zero set of a nonzero real
polynomial, hence a real hypersurface of Lebesgue measure zero in $\R^4$
(Theorem~\ref{thm:a}). The iteration converges to $(\zeta_1,\zeta_2)$ when $|t|<1$ and to
$(\zeta_2,\zeta_1)$ when $|t|>1$.
\item For every $n\ge3$ this fails: there are polynomials with simple zeros for which $E_n$
contains a nonempty open set, hence has positive Lebesgue measure (Theorem~\ref{thm:b}). The
Weierstrass iteration is therefore \emph{not} generally convergent off a measure-zero set for
$n\ge3$.
\end{itemize}

\end{document}
